% =============================================================================
% test-abnt-numeracao.tex
% Documento de teste para o pacote abnt-numeracao.sty
%
% Demonstra todos os recursos do pacote conforme NBR 6024:2012.
% Cada bloco indica qual requisito da norma está sendo exercitado.
%
% Compilação:
%   TEXINPUTS=./src: latexmk -pdf -output-directory=tests tests/test-abnt-numeracao.tex
% =============================================================================

\documentclass[a4paper, 12pt]{report}
\usepackage[T1]{fontenc}
\usepackage[brazil]{babel}
\usepackage{abnt-numeracao}
\usepackage{hyperref}  % checklist item 10: compatibilidade
\usepackage{setspace}
\onehalfspacing

\begin{document}

% =====================================================================
% CAPÍTULO 1 — Exercita: primária em nova página (4), formatação
% hierárquica (2), numeração sem ponto (3), espaçamento (5)
% =====================================================================
\chapter{Introdução}

Este capítulo demonstra a formatação da seção primária: título em maiúsculas
e negrito, alinhado à esquerda, com o indicativo numérico separado por espaço
(sem ponto, hífen ou travessão).

% Checklist 2, 3: formatação hierárquica, numeração sem ponto
\section{Contexto da pesquisa}

A seção secundária aparece em maiúsculas, sem negrito. O indicativo ``1.1''
não tem ponto final.

\subsection{Delimitação do tema}

A seção terciária aparece em negrito, sem maiúsculas automáticas. O indicativo
é ``1.1.1''.

\subsubsection{Aspectos específicos do recorte}

A seção quaternária aparece em fonte normal, sem negrito nem itálico. O
indicativo é ``1.1.1.1''.

\paragraph{Considerações preliminares sobre o recorte temático}

A seção quinária aparece em itálico. O indicativo é ``1.1.1.1.1''. Este
parágrafo demonstra a formatação do quinto nível hierárquico.

% =====================================================================
% CAPÍTULO 2 — Exercita: primária em nova página (4),
% título longo com hanging indent (7),
% \label e \ref em título com hyperref (10)
% =====================================================================
\chapter{Revisão da literatura sobre metodologias de pesquisa em ciências
sociais aplicadas\label{cap:revisao}}

Este capítulo demonstra um título de seção primária longo, que ocupa mais de
uma linha. A segunda linha deve estar alinhada sob a primeira letra do título,
não sob o indicativo numérico.

O capítulo anterior pode ser referenciado: ver Capítulo~\ref{cap:revisao}.

\section{Abordagens qualitativas e suas contribuições para a compreensão de
fenômenos sociais complexos\label{sec:quali}}

Esta seção demonstra um título de seção secundária longo. O indicativo é
``2.1'' e a segunda linha deve alinhar sob a primeira letra do título.

A Seção~\ref{sec:quali} contém referência cruzada no título, testando a
compatibilidade de \verb|\MakeTextUppercase| com \verb|hyperref|.

% =====================================================================
% CAPÍTULO 3 — Exercita: alíneas (8), subalíneas (9)
% =====================================================================
\chapter{Metodologia}

% Checklist 8: alíneas (seção 4.2)
% Texto antes das alíneas termina em dois pontos (seção 4.2.b)
Os procedimentos metodológicos foram organizados nas seguintes etapas:

\begin{alineas}
  \item revisão bibliográfica em bases de dados nacionais e internacionais;
  \item coleta de dados primários por meio de entrevistas semiestruturadas;
  % Checklist 9: subalíneas (seção 4.3)
  \item análise dos dados coletados, incluindo:
    \begin{subalineas}
      \item codificação das entrevistas transcritas;
      \item categorização temática dos conteúdos identificados;
      \item triangulação com dados documentais.
    \end{subalineas}
  \item elaboração do relatório final com os resultados consolidados.
\end{alineas}

As etapas acima foram executadas sequencialmente ao longo de doze meses.

% Demonstrar mais seções para variedade
\section{Participantes}

A amostra foi composta por trinta participantes selecionados por conveniência.

\section{Instrumentos de coleta}

Foram utilizados dois instrumentos principais.

% =====================================================================
% Título sem indicativo numérico, centralizado
% Exercita: checklist item 6 (seção 4.1.h / 14724 5.2.3)
% =====================================================================
\chapter*{Referências}
\addcontentsline{toc}{chapter}{REFERÊNCIAS}

ASSOCIAÇÃO BRASILEIRA DE NORMAS TÉCNICAS. \textbf{NBR 6024}: informação e
documentação: numeração progressiva das seções de um documento: apresentação.
Rio de Janeiro, 2012.

\end{document}
